\section{Restricted kernels: embeddings, locality, and flat residuals}
\label{sec:restricted-kernels}
The full labelled rows of an affine lift carry more information than a
single weighted sum of them, which is a contact kernel. This section
studies what factorizations of \emph{restricted kernels} can and cannot
prove about the number of factors. These are prescribed kernels, such as
the contact kernel of a product of spheres or of projective spaces, used
in place of the full slack operator. We also give a nullity obstruction
that applies when the mixed contact curvature is zero. Each result
states its factorization class. A factorization of a restricted kernel
need not extend to the full slack operator of a convex body.

\subsection{Projective contact forces an embedding}
Write $P_x=xx^T$ for $x\in\Sph^{r-1}$ and let
$\mathbb{RP}^{r-1}$ denote the space of such rank-one projectors. The
contact kernel on this space is
\[
 \kappa([x],[y])=1-\tr(P_xP_y)=1-(x^Ty)^2.
\]
This kernel is \emph{strict-diagonal}: it vanishes exactly on the
diagonal. A factorization is \emph{nowhere-zero} if every selected primal
and dual factor is nonzero at every point, not merely on an open dense
subset.

\begin{theorem}[Projective phase embedding]\label{thm:projective-phase}
Let $r\geq3$, $k\geq1$, and $\lambda_a>0$. Suppose that on
$M=(\mathbb{RP}^{r-1})^k$ the kernel
\[
 \kappa_k(u,v)=\sum_{a=1}^k\lambda_a
                  \bigl(1-(x_a^Ty_a)^2\bigr)
       =\sum_{i=1}^L\ip{A_i(u)}{B_i(v)}
\]
has a globally $C^1$, nowhere-zero factorization with $A_i,B_i\in Q_3$.
Then $M$ has a $C^1$ embedding in $\R^L$. In particular, with
$m=\lceil\log_2 r\rceil$,
\begin{equation}\label{eq:projective-count}
 L\geq\max\{k(r-1)+1,\ k(2^m-1)\}.
\end{equation}
For one projective plane ($r=3$, $k=1$), the minimum of $L$ in this
class is four. For all $r,k$ there is a smooth nowhere-zero factorization with
\begin{equation}\label{eq:projective-upper}
 L=k\left(\frac{r(r+1)}2-2\right).
\end{equation}
\end{theorem}
\begin{proof}
Nonnegative summands and the diagonal zero force complementarity in
every factor pair. Since both factors are nonzero, their normalizations have the
form
\[
 A_i(u)=a_i(u)(1,p_i(u)),\qquad
 B_i(u)=b_i(u)(1,-p_i(u)),\qquad a_i,b_i>0,\quad p_i(u)\in\Sph^1.
\]
The fundamental group of $M$ is finite, so every homomorphism from it to
$\pi_1(\Sph^1)=\mathbb Z$ vanishes. The covering lifting criterion
therefore gives a global $C^1$ angle $\theta_i:M\to\R$ with
$p_i=(\cos\theta_i,\sin\theta_i)$. Equivalently,
$H^1(M;\mathbb Z)=0$; mod-two first cohomology would not suffice.
Taking the negative mixed derivative on the diagonal gives
\[
 2\sum_a\lambda_a g_a
       =\sum_i a_i b_i\,d\theta_i\otimes d\theta_i,
\]
where $g_a$ is the quotient round metric on the $a$th projective space.
Only first derivatives of the factors enter. Since the left side is
positive definite, $\Theta=(\theta_1,\ldots,\theta_L)$ is an immersion. If
$\Theta(u)=\Theta(v)$, every phase agrees and all pairings
$\ip{A_i(u)}{B_i(v)}$ vanish. The strict-diagonal property then gives
$u=v$. Compactness makes this injective immersion an embedding.

For completeness, we include the characteristic-class calculation
behind the numerical bound. Put $n=r-1$ and let $a$ generate
$H^1(\mathbb{RP}^n;\mathbb F_2)$. The tangent bundle is
$\operatorname{Hom}(\gamma,\gamma^\perp)$, where $\gamma$ is the
canonical line bundle, so
$T\mathbb{RP}^n\oplus\varepsilon^1\cong(n+1)\gamma$.
The Whitney product formula \citep{MilnorStasheff1974} gives, for the
normal bundle $\eta$ of an embedding,
\[
 w(\eta)=(1+a)^{-r}
   =\sum_{j=0}^n \binom{n+j}{j}a^j.
\]
The coefficient is odd exactly when the binary expansions of $n$ and
$j$ have no common nonzero bit. This also follows directly by expanding
$(1+a)^{-r}$ over $\mathbb F_2$. The largest such $j\leq n$ is
$j_*=2^m-r$, the complement of the first $m$ bits of $n$.
For the product, $w(\eta)=\prod_a(1+a_a)^{-r}$ has the nonzero monomial
$\prod_a a_a^{j_*}$. Hence the normal rank is at least $kj_*$ and
$L\geq k(2^m-1)$. Independently, a nonempty compact manifold of dimension
$kn$ cannot embed in $\R^{kn}$, by invariance of domain. This proves
\eqref{eq:projective-count}. The projective plane cannot embed in
$\R^3$: a compact embedded hypersurface of Euclidean space is orientable,
whereas $\mathbb{RP}^2$ is not. Thus its lower bound is four.

For the construction put $D=r(r+1)/2-1$ and use the normalized Veronese
map into the unit sphere of traceless symmetric matrices,
\[
 X([x])=\sqrt{\frac r{r-1}}\left(xx^T-\frac Ir\right)
       \in\Sph^{D-1},\qquad
 \kappa(u,v)=\frac{r-1}{r}\bigl(1-\ip{X(u)}{X(v)}\bigr).
\]
Its image is a proper compact subset of that sphere. Choose a pole $p$
outside it, set $c(u)=1-\ip p{X(u)}>0$, and let
$F(u)=(X(u)-\ip p{X(u)}p)/c(u)\in\R^{D-1}$.
The scalar null-paraboloid maps
\begin{equation}\label{eq:scalar-chord-maps}
 U(t)=\left(\frac{1+t^2}2,\frac{1-t^2}2,t\right),\qquad
 V(t)=\left(\frac{1+t^2}2,-\frac{1-t^2}2,-t\right)
\end{equation}
satisfy $\ip{U(t)}{V(w)}=(t-w)^2/2$. The stereographic identity is
\[
 1-\ip{X(u)}{X(v)}
       =\frac{c(u)c(v)}2\norm{F(u)-F(v)}^2.
\]
Thus $\sqrt{(r-1)/r}\,c(u)U(F_j(u))$ and the corresponding $V$ factors
give $D-1$ smooth nowhere-zero channels. Concatenation with scales
$\sqrt{\lambda_a}$ proves \eqref{eq:projective-upper}, and gives four
when $r=3$.
\end{proof}

The bundle obstruction itself is classical. The step specific to
factorizations is the passage from positive complementary Lorentz phases
to an injective Euclidean phase map. The theorem concerns only the
displayed rank-one kernel and its nowhere-zero factorizations. In
particular, its finite upper bound is consistent with Fawzi's theorem
that the full cone
$\mathbb S_+^3$ has no finite second-order-cone lift
\citep{Fawzi2019SOC}.

\subsection{What local tests and separated constructions can show}
Fix $q\geq2$, $k\geq2$, positive weights with $\sum_a\lambda_a=1$, and
\begin{equation}\label{eq:sphere-joint-kernel}
 M=(\Sph^q)^k,\qquad
 K(x,z)=\sum_{a=1}^k\lambda_a(1-x_a^Tz_a),\qquad
 n=kq,\quad D=k(q+1).
\end{equation}
For arbitrary globally $C^1$ factorizations of this joint kernel with
$Q_3$ factors, Proposition~\ref{prop:joint} and the construction after it
give $n+1\leq L\leq n+\lceil k/2\rceil$. The results below show what
several natural attempts to close this gap can and cannot prove.

\begin{proposition}[Exact local saturation]\label{prop:local-saturation}
Every finite collection of primal and dual points of $M$ lies in a product
of punctured-sphere charts on which $K$ has a smooth nowhere-zero
factorization with exactly $n$ channels. The identity holds on the whole
product chart, not merely on the finite samples.
\end{proposition}
\begin{proof}
Choose poles $p_a\in\Sph^q$ outside the finite projections of the sample
sets. On the punctured sphere set
\[
 c_a(x_a)=1-p_a^Tx_a,\qquad
 F_a(x_a)=\frac{x_a-(p_a^Tx_a)p_a}{c_a(x_a)}\in\R^q.
\]
For each of the $q$ coordinates of $F_a$, use the two factors
\[
 \sqrt{\lambda_a}c_a(x_a)U((F_a)_j(x_a)),\qquad
 \sqrt{\lambda_a}c_a(z_a)V((F_a)_j(z_a))
\]
from \eqref{eq:scalar-chord-maps}. The stereographic chord identity gives
exactly the $a$th summand of \eqref{eq:sphere-joint-kernel}. There are
$kq$ channels in total.
\end{proof}
Consequently, neither a finite submatrix of the kernel nor finitely many
of its jets at sample pairs can distinguish the local count $n$ from the
larger global smooth count. This concerns the kernel only; it does not
say whether finitely many samples of
\emph{prescribed factor maps} determine their possible global extension.

\begin{proposition}[Partition separation]\label{prop:partition-channels}
Suppose the source coordinates are partitioned into nonempty groups
$G_1,\ldots,G_m$, and the channel labels into groups $I_1,\ldots,I_m$,
such that both factors of every channel in $I_g$ depend only on the
coordinates in $G_g$. Every global $C^1$ factorization of $K$ of this form
satisfies $L\geq n+m$. Equality is attainable for every partition into
singletons and pairs.
\end{proposition}
\begin{proof}
All diagonal pairings vanish separately. Freeze every group other than
$G_g$ at equal primal and dual arguments. The remaining channels then
factor exactly $\sum_{a\in G_g}\lambda_a(1-x_a^Tz_a)$.
The proof of Proposition~\ref{prop:joint} applies to any product of
simply connected spheres: if its $q|G_g|$-dimensional contact metric were
factored with that many $Q_3$ channels, each channel would have rank one
everywhere; the real phase map would be a local diffeomorphism from a
compact manifold to $\R^{q|G_g|}$. Therefore
$|I_g|\geq q|G_g|+1$. Summing proves the bound. A singleton uses the
$q+1$ scalar chord maps applied to its ordinary coordinates. A pair uses
the $2q+1$ stereographic channels in Proposition~\ref{prop:joint}.
Rescale by the sum of its weights when that sum is not one.
\end{proof}

Another restricted class consists of representations by finitely many
\emph{common-scale squared-distance kernels}:
\begin{equation}\label{eq:chordal-class}
 K(x,z)=\sum_{g=1}^m\frac{r_g(x)r_g(z)}2
                  \norm{F_g(x)-F_g(z)}^2,
 \qquad r_g>0,\quad F_g:M\to\R^{L_g},\quad L=\sum_g L_g,
\end{equation}
with globally $C^1$ maps. Each scalar coordinate produces one $Q_3$
channel by \eqref{eq:scalar-chord-maps}. Different summands may have the
same scale function.

\begin{proposition}[Functional inertia]\label{prop:chordal-inertia}
Let $e$ be the number of summands in \eqref{eq:chordal-class} for which
the image of the feature map $F_g$ is not contained in a Euclidean
sphere. Then
\begin{equation}\label{eq:chordal-inertia}
 L\geq D-e\geq D-m.
\end{equation}
In particular an $L=n+1$ representation in this class requires $e\geq k-1$.
Within the narrower class in which each summand is itself a positive
block-additive kernel
\begin{equation}\label{eq:positive-block-chord}
 \frac{r_g(x)r_g(z)}2\norm{F_g(x)-F_g(z)}^2
   =\sum_{a\in G_g}\mu_{ga}(1-x_a^Tz_a),\qquad
 \mu_{ga}>0,\quad \sum_{g:a\in G_g}\mu_{ga}=\lambda_a,
\end{equation}
the exact minimum count is $n+\lceil k/2\rceil$, even when the groups
overlap.
\end{proposition}
\begin{proof}
Set $\alpha_g=r_g\norm{F_g}^2/2$, $\beta_g=r_g$, and
$h_{gj}=r_g(F_g)_j$. The corresponding kernel is
\[
 \alpha_g(x)\beta_g(z)+\beta_g(x)\alpha_g(z)
                  -\sum_jh_{gj}(x)h_{gj}(z).
\]
Its coefficient form has at most $L_g+1$ negative directions. If the
image lies on a sphere, translate the center to zero without changing
squared distances. Then $\alpha_g$ is a nonnegative constant multiple
of $\beta_g$ and the bound improves to $L_g$ negative directions.
Any pullback to the actual span of functions has no larger negative
index. Conversely, the constant function and the $D$ coordinate
functions of $X(x)=(\sqrt{\lambda_a}x_a)_a$ are linearly independent,
and $K(x,z)=1-X(x)^TX(z)$. Choose finitely many evaluation points whose
feature evaluation matrix has full column rank. Their kernel matrix has
exactly $D$ negative eigenvalues. Applying the preceding upper bounds to
this same finite matrix proves \eqref{eq:chordal-inertia}.

For \eqref{eq:positive-block-chord}, a nonsingleton group contributes
at least $|G_g|(q+1)-1$ channels by the same inertia argument. A singleton
contributes at least $q+1$ by the compact phase obstruction. If
$I=\sum_g|G_g|$ and $t$ groups are nonsingletons, then $I\geq k$,
$t\leq\lfloor I/2\rfloor$, and
\[
 L\geq(q+1)I-t\geq(q+1)I-\lfloor I/2\rfloor
       \geq(q+1)k-\lfloor k/2\rfloor.
\]
Disjoint pairs and, if needed, one singleton attain this value.
\end{proof}
In the paired construction, the image of each stereographic pair lies
in no Euclidean sphere. Indeed, the inverse image of a Euclidean sphere under
stereographic projection is contained in a hyperplane section of the
ambient unit sphere, whereas a weighted product of two full spheres
affinely spans its entire ambient space. The image of a singleton lies in a sphere. Thus this construction also
attains \eqref{eq:chordal-inertia}, with $e=\lfloor k/2\rfloor$.
In the general class \eqref{eq:chordal-class}, $m$ and $e$ can be
larger; the paired construction is proved optimal only in the class
\eqref{eq:positive-block-chord}.

\begin{example}[A channel can depend on every source]\label{ex:all-block-channel}
Choose a pole $p=(p_a)_a\in\Sph^{D-1}$ with every $p_a\ne0$ and
$(\norm{p_a})_a\ne(\sqrt{\lambda_a})_a$. Then
\[
 \max_{x\in M}p^TX(x)=\sum_a\sqrt{\lambda_a}\norm{p_a}<1.
\]
Global stereographic projection of $X(M)$ from $p$ therefore gives
$D-1$ smooth nowhere-zero channels. To verify dependence on all sources,
choose a unit $v\in p^\perp$ such that each $v_a\ne0$ and
$v_a\perp p_a$. This is possible because each block has dimension at
least three. One coordinate is $f=v^TX/c$, where $c=1-p^TX$.
The time component of its primal factor is
$T=c(1+f^2)/2$. Fix a source $a$. On an open subset of its sphere there
are tangent variations with $d(v^TX)=0$ and $dc\ne0$, because $v_a$ and
$p_a$ are independent. At points with $f\ne\pm1$ these give
$dT=(1-f^2)dc/2\ne0$. Such points exist: the equation $v^TX=\pm c$
cannot hold on an open subset of the product, since its two sides would
give a nontrivial affine relation among the constant and coordinate
functions. This argument applies to every $a$. Thus one channel can
depend nontrivially on all $k$ sources. Positivity alone therefore
does not let a proof for unrestricted factorizations assume that each
channel depends on one source or one pair of sources.
\end{example}

\subsection{Column ranks detect residuals with zero contact curvature}
The next theorem uses finite-dimensional linear algebra instead of
contact differentiation, and it needs no topology. For a kernel
$F_a:X_a\times X_a\to\R$, its \emph{column rank} is the dimension of the
span of the functions $F_a(\,\cdot\,,v)$ on $X_a$.

\begin{theorem}[A column-rank nullity criterion]\label{thm:rado-nullity}
Let $X_1,\ldots,X_b$ be nonempty sets and $M=\prod_aX_a$.
Suppose nonnegative kernels $F_a$ vanish on the diagonal, have finite
nonzero column ranks $\rho_a$, and admit the labelled real-PSD
factorization
\begin{equation}\label{eq:residual-factors}
 F_a(x_a,v)=\sum_{i=1}^L\tr\bigl(A_i(x)B_i^a(v)\bigr),
 \qquad A_i(x),B_i^a(v)\in\mathbb S_+^{r_i},\quad r_i\leq R.
\end{equation}
If, for every nonempty $J\subseteq\{1,\ldots,b\}$,
\begin{equation}\label{eq:rado-criterion}
 1+\sum_{a\in J}(\rho_a-1)>
                   \frac{R+1}{2}(|J|-1),
\end{equation}
then some $x\in M$ satisfies
$\sum_i\operatorname{nullity}A_i(x)\geq b$.
A sufficient uniform condition, when $\rho_a\geq\rho>0$, is
\begin{equation}\label{eq:rado-uniform}
 (R+3)(b-1)<2b\rho.
\end{equation}
For $b=1$ both conditions hold for every $R$.
\end{theorem}
\begin{proof}
Assume $\sum_i\operatorname{nullity}A_i(x)\leq b-1$ for every $x\in M$.
In
$\mathcal H=\bigoplus_i\R^{r_i}$ write $A=\bigoplus_iA_i$ and
$B^a=\bigoplus_iB_i^a$. Diagonal contact implies
$\operatorname{Ran}B^a(x_a)\subseteq\ker A(x)$. For each $a$ let
$S_a$ be the union of the nonzero vectors in the ranges of all $B^a(v)$.
It is nonempty because $\rho_a>0$. If there were linearly independent
vectors $w_a\in S_a$, then choosing the corresponding $v_a$ and setting
$x=(v_a)_a$ would put $b$ independent vectors in $\ker A(x)$.

The independent-transversal criterion of Rado \citep{Rado1942} therefore
gives a nonempty $J$, with $j=|J|$, such that
\[
 \dim\spanop\bigcup_{a\in J}S_a\leq j-1.
\]
We give a direct proof that also covers infinite sets $S_a$. For
finite-dimensional subspaces $V_a=\spanop S_a$, independent vectors
$v_a\in V_a$ exist if $\dim\sum_{a\in J}V_a\geq|J|$ for every $J$.
Prove this by induction. If a nonempty proper subset is tight, first
choose its independent vectors and then apply induction to the remaining
subspaces in the quotient by its sum. Otherwise every nonempty proper
subset has dimension at least its cardinality plus one; select a nonzero
vector in the last subspace and quotient by its span, preserving the
inductive inequalities for the others. Finally, multilinearity of the
exterior product expands independent choices from the $V_a$ into at least
one independent choice from the generating sets $S_a$. This proves the
claimed criterion without a finiteness assumption on those sets.

Put $V_i=\spanop_{a\in J,v}\operatorname{Ran}B_i^a(v)$ and
$m_i=\dim V_i$. The range of a block-diagonal matrix contains each
block-supported range vector. Consequently the span of $\bigcup_{a\in J}S_a$ is exactly
$\bigoplus_i V_i$, not merely a subspace of it, and
$\sum_i m_i\leq j-1$. All the dual matrices for $a\in J$ lie in
$\bigoplus_i\operatorname{Sym}(V_i)$, whose dimension is at most
\[
 \sum_i\frac{m_i(m_i+1)}2
       \leq\frac{R+1}2\sum_i m_i
       \leq\frac{R+1}2(j-1).
\]
Pairing these matrices with $A(x)$ bounds by the same number the span of
all cylindrical column functions $x\mapsto F_a(x_a,v)$, $a\in J$.
On the other hand, a relation $\sum_{a\in J}f_a(x_a)=0$ makes each
$f_a$ constant, by varying one coordinate at a time. Its relation space
has dimension at most $j-1$, so the cylindrical span has dimension at
least $\sum_{a\in J}\rho_a-(j-1)$. This contradicts
\eqref{eq:rado-criterion}.
For the uniform assertion, with $\rho$ in place of each $\rho_a$,
twice the difference of the two sides of \eqref{eq:rado-criterion} is
$R+3+j(2\rho-R-3)$, an affine function of $1\leq j\leq b$.
It is positive at $j=1$, and at $j=b$ positivity is exactly
\eqref{eq:rado-uniform}. For $b=1$ this condition is automatic.
\end{proof}

\begin{corollary}[Antipodal separation and squared slack]\label{cor:flat-nullity}
Let $s\geq3$, $X_a=\Sph^{s-1}$, and assume the kernels in
\eqref{eq:residual-factors} are continuous. If each $F_a(u,u)=0$ and
$F_a(u,\tau_a(u))>0$, where $\tau_a$ is topologically conjugate to the
antipodal map, then the nullity conclusion of
Theorem~\ref{thm:rado-nullity} holds whenever
\begin{equation}\label{eq:orientation-cap}
 (R+3)(b-1)<2bs.
\end{equation}
For the squared kernels
$F_a(u,v)=(1-u^Tv)^2/2$, the column rank is exactly
$\rho=s(s+3)/2$. Thus \eqref{eq:orientation-cap} can be replaced by
the less restrictive sufficient condition
\begin{equation}\label{eq:squared-cap}
 (R+3)(b-1)<b s(s+3).
\end{equation}
Without any cap condition, every factorization of these squared kernels
satisfies
\begin{equation}\label{eq:squared-dimension}
 \sum_i\frac{r_i(r_i+1)}2\geq1+b\left(\frac{s(s+3)}2-1\right).
\end{equation}
\end{corollary}
\begin{proof}
Choose a basis $f_1,\ldots,f_{\rho_a}$ of the column space and write
$F_a(u,v)=\sum_\ell f_\ell(u)c_\ell(v)$. The coefficient map
$c:\Sph^{s-1}\to\R^{\rho_a}$ is continuous: select $\rho_a$
evaluation points for which the basis evaluation matrix is invertible
and recover $c$ from those evaluations. If $c(v)=c(\tau_a(v))$, evaluation
at $u=v$ contradicts the assumed separation. After conjugating $\tau_a$
to the antipodal map, the Borsuk--Ulam theorem
\citep[Section 2.2]{Hatcher2002} implies $\rho_a\geq s$.
Theorem~\ref{thm:rado-nullity} proves \eqref{eq:orientation-cap}.

For the squared kernel the odd parts of the columns span all $s$ linear
coordinate functions. The even parts span all quadratic forms restricted
to the sphere: their coefficient matrices are proportional to $I+vv^T$.
A matrix $C$ orthogonal to every $I+vv^T$ obeys
$\tr C+v^TCv=0$ for every unit $v$, which first makes $C$ scalar and then
zero. The restriction of homogeneous quadratic forms to the sphere is
injective, and the odd and even spaces are independent. Thus
$\rho=s+s(s+1)/2$. This gives \eqref{eq:squared-cap}.
The constant function belongs to the column span, for example by averaging
$F_a(\,\cdot\,,v)$ over the sphere. Coordinate column spaces are
independent modulo this common one-dimensional space, so the labelled
cylindrical column rank is exactly $1+b(\rho-1)$. Pairing with the full
matrix space $\bigoplus_i\mathbb S^{r_i}$ proves
\eqref{eq:squared-dimension}.
\end{proof}

These squared kernels have zero mixed contact curvature, so the contact
Hessian estimate cannot recover this lower bound. If the
selected primal tuples also lie in the closure of a relative-Slater
affine slice, Lemma~\ref{lem:boundary-order} converts the nullity
conclusion into $\nustd\geq b$ for the standard barrier of that slice.

For example, the identity
\[
 1-u^Tv=\frac{1-(u^Tv)^2}{2}+\frac{(1-u^Tv)^2}{2}
\]
splits the ball slack into a projective term and a flat residual. The first term has the
order-$s$ factorization $\tr((I-uu^T)vv^T/2)$. A construction that uses
one such separate block for each source has $b$ primal nullities at every
contact. If the residual uses blocks of order at most $s$,
Corollary~\ref{cor:flat-nullity} forces at least another $b$ nullities at
some simultaneous contact. Moreover, with
$d_s=s(s+1)/2$, \eqref{eq:squared-dimension} gives
\[
 L_{\rm residual}\geq b+
       \left\lceil\frac{b(s-1)+1}{d_s}\right\rceil>b.
\]
Thus this prescribed positive splitting needs at least $2b+1$ factors,
whereas the grouped Schur construction for the unsplit full slacks uses
$2b$. The rank theorem for unrestricted full slacks does not assume such a
splitting. The example shows that a projective summand attaining
equality in the curvature bound need not be an economical part of a
complete formulation.

Lower bounds on PSD factorization size from ordinary matrix rank are
standard \citep{FGPRT2015}. Theorem~\ref{thm:rado-nullity} adds the
combination of the independent-transversal criterion with simultaneous
contact and the blockwise bound on matrix-space dimension. Stating the
criterion with the actual column ranks keeps information that the weaker
antipodal estimate $\rho_a\geq s$ loses. The conclusions concern a
simultaneous primal nullity for a prescribed labelled kernel. Without
further hypotheses, they bound neither the nullity at every contact nor
an arbitrary barrier.
